% acknowledgements.tex
%
% Aetf <aetf@unlimitedcodeworks.xyz>
% Copyright 2016 Aetf <aetf@unlimitedcodeworks.xyz>
%
% multiple1902 <multiple1902@gmail.com>
% Copyright 2011~2012, multiple1902 (Weisi Dai)
%
% Project Home: https://github.com/Aetf/xjtuthesis
%
% It is strongly recommended that you read documentations located at
%   https://github.com/Aetf/xjtuthesis/wiki
% in advance of your compilation if you have not read them before.
%
% This work may be distributed and/or modified under the
% conditions of the LaTeX Project Public License, either version 1.3
% of this license or (at your option) any later version.
% The latest version of this license is in
%   http://www.latex-project.org/lppl.txt
% and version 1.3 or later is part of all distributions of LaTeX
% version 2005/12/01 or later.
%
% This work has the LPPL maintenance status `maintained'.
%
% The Current Maintainer of this work is Aetf.
%

作为少年班的学生，我已在兴庆校区度过了五载春秋。
今朝临近毕业，蓦然回首，方觉时光飞逝。
在这段漫长的旅途中，我似乎不曾留下什么浓墨重彩的印记，
更多时候，只是一个人在迷茫与摸索中缓慢前行。

本篇论文亦是在这般心境与状态下勉力完成。在毕设起步之初，
我深感自身的人工智能基础尚不扎实。因而在整个毕业设计的期间，
我经历了无数次的试错与反复，一边学习相关理论基础，一边在实践中摸爬滚打，
方使这项工作得以艰难而顺利地推进下去。

此项工作的顺利完成，首先要由衷感谢我的毕设指导老师蔺杰老师，
以及同组学长们给予的无私帮助与支持。
同时，也要向开源社区 SMAC 和 MAIC 的开发者们致敬，
他们优秀的实验框架为本文奠定了坚实的基础。

大学生活中，我得到了太多温暖的托举与陪伴。
感谢我深爱的本科社团——团委挑战网，那里是我大学记忆中最闪亮、最温暖的一部分；
感谢李昊老师和夏泽学长，是他们带我初窥科研门径，领略了学术的魅力；
感谢实习期间各位同事悉心的技术指导与工程实践上的关怀；
感谢“交大门”论坛上那些与我一样曾陷于迷茫并给予彼此温情鼓励的朋友们。
同时，感谢父母给予我自由探索的机会，也感谢求学路上遇到的所有老师与同窗。

回望这些年走过的路，我深感自己是幸运的。在交大，我邂逅了许多挚友和才华横溢的同行者。
虽然前进的脚步偶有迟缓，但我相信，那个在迷茫中从未放弃、执着摸索的自己，配得上一个更好的未来。
愿过去的积累化作风起之时的羽翼。祝自己轻舟已过万重山，在未来的广阔天地里，越飞越高。